% =====================================================================
%  Layup Equilibria — 15-minute presentation
%  Compile: pdflatex -interaction=nonstopmode main.tex
%  All slide text is verbatim or near-verbatim from ../main.tex.
% =====================================================================
\documentclass[xcolor=dvipsnames,11pt,table,aspectratio=169]{beamer}

\definecolor{DartmouthGreen}{RGB}{0,105,62}
\setbeamercolor{structure}{fg=DartmouthGreen}
\setbeamercolor{title}{fg=DartmouthGreen}
\setbeamercolor{frametitle}{fg=DartmouthGreen}
\setbeamertemplate{navigation symbols}{}

\usepackage{booktabs}
\usepackage{graphicx}
\usepackage{amsmath, amssymb, bm}
\usepackage{array}
\graphicspath{{./Figures/}}

\begin{document}

% --- Slide 1: Title ---------------------------------------------------
\begin{frame}
\begin{center}
\vspace{0.2cm}
{\Large \textbf{Layup Equilibria}}\\[4pt]
{\large \textbf{An Evolutionary Game-Theoretic Model of\\
Strategic Course Selection at Dartmouth}}\\[14pt]
\small Jack Chou \quad Alejandro Araujo \quad Nicholas Booth \quad Godwin Kangor\\[10pt]
\footnotesize\textcolor{DartmouthGreen}{\textbf{MATH/QSS 30.04: Evolutionary Game Theory}}\\[1pt]
\footnotesize Dartmouth College, May 2026
\end{center}
\end{frame}

% --- Slide 2: Motivation ----------------------------------------------
\begin{frame}
\frametitle{What is a layup?}
\small
\begin{itemize}
\item At Dartmouth, a ``layup'' is a course taken mainly for its
      high expected grade rather than its content.
\item Academic performance is benchmarked against peers through graded
      distributions, class rankings, and competitive post-graduate
      admissions, so each student's best portfolio depends on what
      classmates are doing.
\item A student who picks a rigorous course while peers enroll in easy
      ones takes a double hit: harder material and a harsher relative
      curve.
\item \textbf{What looks like a private study decision is in fact
      strategic.}
\end{itemize}
\end{frame}

% --- Slide 4: The three-type replicator model (condensed) ------------
\begin{frame}
\frametitle{The three-type replicator model}
\small
Partition the student population into three behavioral types:
\begin{itemize}
\itemsep1pt
\item \emph{Purists} (type $A$): few layups, $s_A = 0.10$
\item \emph{Balancers} (type $B$): moderate, $s_B = 0.35$
\item \emph{Strategists} (type $C$): maximize GPA-padding, $s_C = 0.70$
\end{itemize}
\vskip 4pt
\textbf{Payoffs.} Taking layups pays off when you out-pad the crowd, but
each layup's value erodes as more peers pile in. The strength of that
crowding is the competition intensity $\gamma$, weighed against how much
GPA matters ($w$). The population then evolves by replicator dynamics:
types earning above the average grow, those below shrink.
\vskip 6pt
\textbf{Key result.} Once GPA matters enough --- past a threshold of
about 2.6 on the 1--5 scale --- strategists become unbeatable. No mix of
purists or balancers can displace them.
\end{frame}

% --- Slide 6: Phase portraits -----------------------------------------
\begin{frame}
\frametitle{Phase portraits at $w \in \{1, 3, 5\}$}
\begin{center}
\includegraphics[width=0.31\textwidth]{simplex_w1.png}\hfill
\includegraphics[width=0.31\textwidth]{simplex_w3.png}\hfill
\includegraphics[width=0.31\textwidth]{simplex_w5.png}
\end{center}
\vskip 4pt
\scriptsize{
When GPA matters little ($w = 1$) no corner wins and the population
settles between the purist and balancer corners. When it matters more
($w = 3, 5$) everything flows to the strategist corner. As GPA pressure
rises, the population swings from the purist corner toward the
strategist corner.
}
\end{frame}

% --- Slide 7: Four testable predictions -------------------------------
\begin{frame}
\frametitle{Four testable predictions}
\small
The model says more GPA pressure should mean more layups. That gives
four things we can check in the survey:
\begin{itemize}
\itemsep3pt
\item[\textbf{P1}] Students who care more about GPA take more layups.
\item[\textbf{P2}] Career tracks line up by how much they care about GPA
      (pre-medical $>$ finance/consulting $>$ tech $>$ undecided).
\item[\textbf{P3}] Students who dropped or avoided a course to protect
      their GPA take more layups and care more about GPA.
\item[\textbf{P4}] Students who feel more peer pressure toward easy
      courses take more layups themselves.
\end{itemize}
\end{frame}

% --- Slide 8: Survey and sample ---------------------------------------
\begin{frame}
\frametitle{Survey \& sample}
\begin{columns}[T,onlytextwidth]
\begin{column}{0.42\textwidth}
\footnotesize
\begin{itemize}
\itemsep2pt
\item Google Form survey, Spring 2026, $n = 51$ Dartmouth
      undergraduates, distributed through peer networks.
\item Sample skews toward Finance/consulting ($n = 25$, 49\%) and
      Economics majors ($n = 14$, 27\%).
\item Average GPA-importance rating is 3.6 on a 1--5 scale.
\item Most students take layups in 0--40\% of their courses, with a
      smaller tail above 50\%.
\item Average layup share by track ranges from 15\% (graduate/PhD)
      to 35\% (finance/consulting).
\end{itemize}
\end{column}
\begin{column}{0.56\textwidth}
\begin{center}
\includegraphics[width=\linewidth,keepaspectratio]{desc_overview.png}
\end{center}
\end{column}
\end{columns}
\end{frame}

% --- Slide 9: P1 + P2 results -----------------------------------------
\begin{frame}
\frametitle{P1 \& P2: GPA importance and career track}
\begin{columns}[T]
\begin{column}{0.5\textwidth}
\small
\textbf{P1 --- does caring about GPA mean more layups?}
\begin{itemize}
\itemsep2pt
\item Yes, in the predicted direction --- but the effect is tiny and
      not statistically reliable.
\item Caring more about GPA barely changes how many layups a student takes.
\end{itemize}
\vskip 6pt
\textbf{P2 --- do career tracks line up as the model predicts?}
\begin{itemize}
\itemsep2pt
\item No. Tracks differ only weakly, and not in the predicted order.
\item Finance/consulting takes the most layups (35\% of courses), while
      the supposedly ``most strategic'' tracks (pre-medical, graduate/PhD)
      sit near the bottom.
\end{itemize}
\end{column}
\begin{column}{0.5\textwidth}
\begin{center}
\includegraphics[height=4.8cm]{p2_box_s_by_track.png}
\end{center}
\end{column}
\end{columns}
\end{frame}

% --- Slide 10: P3 + P4 results ----------------------------------------
\begin{frame}
\frametitle{P3 \& P4: avoidance and peer pressure}
\begin{columns}[T]
\begin{column}{0.5\textwidth}
\small
\textbf{P3 --- do GPA-protective course-droppers take more layups?}
\begin{itemize}
\itemsep2pt
\item Students who have dropped or avoided a course to protect their GPA
      ($n = 25$) take somewhat more layups, and care more about GPA, than
      those who never have ($n = 26$).
\item Both gaps point the predicted way but stay weak --- the layup gap
      is borderline, the GPA-importance gap is not significant.
\end{itemize}
\vskip 6pt
\textbf{P4 --- does feeling peer pressure mean more layups?}
\begin{itemize}
\itemsep2pt
\item A weak positive link: students who feel the environment pushes
      toward easy courses take slightly more layups.
\item Not strong enough to be statistically conclusive.
\end{itemize}
\end{column}
\begin{column}{0.5\textwidth}
\begin{center}
\includegraphics[height=4.6cm]{p4_scatter_q9_s.png}\\[2pt]
\scriptsize{Each dot is a student: perceived pressure toward easy courses
(left to right) vs.\ their layup share (up).}
\end{center}
\end{column}
\end{columns}
\end{frame}

% --- Slide 10b: Predictions scorecard ---------------------------------
\begin{frame}
\frametitle{Scorecard: how the four predictions fared}
\small
\begin{center}
\renewcommand{\arraystretch}{1.5}
\begin{tabular}{p{5.0cm}p{4.4cm}l}
\toprule
Prediction & What we found & Holds? \\
\midrule
P1: more GPA focus $\to$ more layups & right direction, tiny effect & weakly \\
P2: tracks ordered by GPA focus & wrong order, weak differences & no \\
P3: GPA-droppers take more layups (and care more about GPA) & right direction, borderline & partly \\
P4: peer pressure $\to$ more layups & right direction, weak & weakly \\
\bottomrule
\end{tabular}
\end{center}
\vskip 6pt
\begin{itemize}
\itemsep2pt
\item All four point the predicted way, but none is strong enough to be
      statistically convincing.
\item The biggest miss: most students cluster among the balancers, not
      the strategists the model predicts.
\end{itemize}
\end{frame}

% --- Slide 11: Empirical simplex anomaly ------------------------------
\begin{frame}
\frametitle{Empirical simplex: the central anomaly}
\begin{columns}[T]
\begin{column}{0.45\textwidth}
\scriptsize
Each dot is a career track, placed by its mix of the three types. The
three corners are all-purist, all-balancer, and all-strategist.
\begin{itemize}
\itemsep2pt
\item No track lands in the strategist corner.
\item Only Finance/consulting, Tech/engineering, and Undecided/other have
      any strategists at all (16\%, 22\%, 14\%). Finance/consulting is
      mostly balancers (76\%).
\item The other tracks sit between the purist and balancer corners.
\item \textbf{The model says everyone should pile into the strategist
      corner. The data show the opposite ---} every track sits in the
      balancer region instead.
\end{itemize}
\end{column}
\begin{column}{0.55\textwidth}
\begin{center}
\includegraphics[height=4.8cm]{empirical_simplex.png}
\end{center}
\end{column}
\end{columns}
\end{frame}

% --- Slide 12: Discussion: two anomalies ------------------------------
\begin{frame}
\frametitle{Discussion: two anomalies}
\small
\textbf{The pre-medical reversal.} Pre-meds care about GPA the most
(4.3 out of 5) yet take among the fewest layups (about 17\% of courses,
and only 3 respondents, so read with caution). A likely explanation:
they protect their GPA a different way --- avoiding hard courses
altogether rather than loading up on easy ones. The right lens for them
may be course difficulty, not layup count.\\[6pt]
\textbf{The balancer puzzle.} Most students care about GPA enough that
the model says strategists should take over --- yet the population sits
firmly among the balancers. Three explanations, likely working together:
\begin{itemize}
\itemsep2pt
\item Competition may be far fiercer than our calibration assumes.
\item The model ignores the costs of being an extreme strategist
      (reputation with faculty, advising limits, wanting to actually learn).
\item Our sample over-represents Economics and Finance, whose norms may
      hold back extreme layup-taking.
\end{itemize}
\end{frame}

% --- Slide 13: Policy counterfactual ----------------------------------
\begin{frame}
\frametitle{Policy counterfactual: eliminating enforced medians}
\begin{columns}[T]
\begin{column}{0.55\textwidth}
\small
\begin{itemize}
\itemsep3pt
\item Removing enforced grading medians weakens the link between your
      grade and what classmates do --- in the model, it halves the
      competitive pressure.
\item The layup share the model pushes toward then \emph{doubles} at
      every level of GPA focus (up to a 100\% cap).
\item \textbf{The implication is troubling.} Removing medians pushes the
      layup share \emph{up}, especially for the low-to-moderate GPA-focus
      students reformers most hope to steer toward harder courses.
\end{itemize}
\end{column}
\begin{column}{0.45\textwidth}
\centering
\footnotesize
\begin{tabular}{cccc}
\toprule
GPA focus & Target now & Medians removed & Change \\
\midrule
1 & 17\% & 33\%  & $+17$ \\
2 & 50\% & 100\% & $+50$ \\
3 & 83\% & 100\% & $+17$ \\
4 & 100\% & 100\% & $0$ \\
5 & 100\% & 100\% & $0$ \\
\bottomrule
\end{tabular}\\[4pt]
\scriptsize Competitive pressure halved, target capped at 100\%.
\end{column}
\end{columns}
\end{frame}

% --- Slide 14: Limitations & future directions ------------------------
\begin{frame}
\frametitle{Limitations \& future directions}
\scriptsize
\textbf{Limitations:}
\begin{itemize}
\itemsep1pt
\item The sample ($n = 51$) is a single-institution convenience sample
      drawn through the authors' peer networks, with no random
      selection and all key variables self-reported.
\item The resulting skew toward Finance/consulting and Economics is
      the most visible expression of this design limitation.
\item The model's competition and baseline-value parameters are set by
      hand, not estimated from data.
\end{itemize}
\vskip 4pt
\textbf{Future work should pursue:}
\begin{itemize}
\itemsep1pt
\item Transcript-level data, so the competition parameter can be
      estimated rather than assumed.
\item Multi-institution comparisons --- Princeton's 2004 grade-deflation
      policy offers a natural experiment.
\item A richer model that adds course difficulty, not just layup count.
\end{itemize}
\end{frame}

% --- Slide 15: Three contributions / takeaway -------------------------
\begin{frame}
\frametitle{Takeaways}
\small
\begin{enumerate}
\itemsep4pt
\item \textbf{A three-type model} (purists, balancers, strategists)
      where layups pay off until everyone piles in. It predicts a sharp
      tipping point: once GPA matters enough, strategists take over.
\item \textbf{Four predictions tested} against $n = 51$ Dartmouth
      undergraduates. All four lean the predicted way but mostly fall
      short of statistical significance. Two big surprises stand out ---
      the pre-medical reversal and the population clustering among the
      balancers, not the strategists.
\item \textbf{A median-elimination counterfactual} that raises ---
      not lowers --- the equilibrium layup ratio for students with
      low-to-moderate GPA importance.
\end{enumerate}
\end{frame}

% --- Slide 16: Acknowledgements (placeholder) -------------------------
\begin{frame}
\begin{center}
\vskip 1cm
{\Large \textbf{\textcolor{DartmouthGreen}{Thank you.}}}\\[14pt]
\normalsize Questions?\\[18pt]
\textcolor{red}{\textbf{[PLACEHOLDER: acknowledgements / advisor / data
provenance note — fill in before delivery]}}
\end{center}
\end{frame}

\end{document}
